\documentclass[11pt,a4paper]{article}

\usepackage[a4paper,margin=2.5cm]{geometry}
\usepackage{amsmath}
\usepackage{amssymb}
\usepackage{hyperref}

\title{\textbf{Topological Integrity Gravity (TIG):\\
A Regular Black Hole Solution}}

\author{Kai}

\date{\today}

\begin{document}

\maketitle

\begin{abstract}
We present a regular black hole solution within the framework of Topological Integrity Gravity (TIG). By introducing an effective integrity scale, the classical Schwarzschild singularity at \( r \to 0 \) is replaced by a finite core. The resulting metric is asymptotically consistent with General Relativity and avoids divergent curvature behavior in the central region.
\end{abstract}

\section{Introduction}

In classical General Relativity, the Schwarzschild solution describes the exterior geometry of a static, spherically symmetric mass. However, the classical solution contains a curvature singularity at \( r = 0 \), where the geometric description breaks down.

Topological Integrity Gravity introduces an intrinsic stabilization scale that regularizes the high-curvature regime and replaces the singularity with a finite core.

\section{Metric Ansatz}

We consider a static, spherically symmetric spacetime of the form:
\begin{equation}
ds^2 = -A(r)\,dt^2 + \frac{dr^2}{A(r)} + r^2 d\Omega^2 .
\end{equation}

\section{Regularized Metric Function}

The proposed regular metric function is:
\begin{equation}
A(r) = 1 - \frac{2GM r^2}{r^3 + r_0^3},
\end{equation}
where \( M \) is the total mass, \( G \) is Newton's gravitational constant, and \( r_0 \) denotes the integrity scale.

\section{Asymptotic Behavior}

For \( r \gg r_0 \), one obtains:
\begin{equation}
A(r) \approx 1 - \frac{2GM}{r}.
\end{equation}

Thus, the classical Schwarzschild solution is recovered in the large-distance limit.

\section{Near-Core Behavior}

For \( r \to 0 \), the metric function behaves as:
\begin{equation}
A(r) \approx 1 - \frac{2GM}{r_0^3} r^2 .
\end{equation}

This corresponds to a de Sitter-like core with an effective cosmological constant:
\begin{equation}
\Lambda_{\mathrm{eff}} = \frac{6GM}{r_0^3}.
\end{equation}

\section{Regularity}

Because the metric function remains finite at \( r = 0 \), the central singularity of the Schwarzschild solution is removed. The core behaves as a finite-curvature region rather than a divergent point.

\section{Physical Interpretation}

The parameter \( r_0 \) represents an effective minimal length scale associated with the integrity structure of spacetime. It may be interpreted as a geometric ultraviolet cutoff that prevents collapse into a singular state.

\section{Conclusion}

The proposed TIG black hole solution preserves the Schwarzschild limit at large distances while replacing the central singularity with a finite de Sitter-like core. This provides a first concrete application of the TIG framework to high-curvature gravitational systems.

\end{document}
